\begin{table}

\caption{\label{tab:evaluation_consistency_check_dispersion}Dispersion of the per-draw CLT SE $\hat\sigma_{M_e}^{(k)}/\sqrt{M_e}$ across the 100 draws underlying Table~\ref{tab:evaluation_consistency_check}. Even though the RMS-combined SE there is well-calibrated on average, a single real run's own SE is highly unstable, especially at higher dimension: at $d=25$ it can be off by a factor of 16 depending on which evaluation sample was drawn. Median/RMS $<1$ at every cell shows this instability is right-skewed, not just wide: a typical single run understates the RMS-combined SE (down to $\approx$0.6 at $d=25$), while only a minority of runs that catch a rare large-jump path overstate it -- this same skew is what produces the asymmetric multipliers in Table~\ref{tab:evaluation_consistency_check_empirical_ci}. Coverage is the practical payoff of that instability: the fraction of the 100 draws' own nominal 95\% CLT intervals ($\tilde V_0 \pm 1.96 \cdot \hat\sigma_{M_e}/\sqrt{M_e}$) that contain the empirical mean across all 100 draws; with only 100 draws its own standard error is $\approx$2.2\%, so the values shown are not meaningfully distinguishable from each other or from nominal 95\%, but do show the instability above does not translate into badly broken coverage.}
\footnotesize
\setlength{\tabcolsep}{4pt}
\centering
\begin{tabular}[t]{clcccccc}
\toprule
Dim & Alg & SE min & SE max & Max/Min & CV (\%) & Median/RMS & Coverage (\%)\\
\midrule
1 & Laguerre (deg 2) & 191.4 & 1159.0 & 6.1 & 48 & 0.79 & 94\\
1 & RNN & 188.7 & 1158.7 & 6.1 & 49 & 0.79 & 88\\
\addlinespace
10 & Laguerre (deg 2) & 538.2 & 2411.4 & 4.5 & 36 & 0.84 & 94\\
10 & RNN & 538.2 & 2411.2 & 4.5 & 36 & 0.84 & 94\\
\addlinespace
25 & Laguerre (deg 2) & 872.0 & 13991.6 & 16.0 & 95 & 0.59 & 91\\
25 & RNN & 877.6 & 13991.9 & 15.9 & 94 & 0.61 & 92\\
\bottomrule
\end{tabular}
\end{table}
